% BHU PYQ -- Manually transcribed & error-corrected
% Paper: GLB-502: Igneous Petrology, Mineralogy and Crystallography
% Exam: B.Sc. (Hons.) Semester V Examination, 2013-14

\documentclass[11pt,a4paper]{article}
\usepackage[a4paper,top=2.5cm,bottom=2.5cm,left=2.8cm,right=2.8cm,headheight=14pt]{geometry}
\usepackage{amsmath,amssymb}
\usepackage[T1]{fontenc}
\usepackage[utf8]{inputenc}
\usepackage{lmodern,microtype,fancyhdr,enumitem,hyperref,lastpage,parskip}

\hypersetup{colorlinks=true,urlcolor=black,linkcolor=black,
  pdftitle={GLB-502 Igneous Petrology, Mineralogy and Crystallography -- BHU B.Sc. Sem V 2013-14}}
\pagestyle{fancy}\fancyhf{}
\renewcommand{\headrulewidth}{0.4pt}\renewcommand{\footrulewidth}{0.4pt}
\fancyhead[L]{\small   \textit{Banaras Hindu University}}
\fancyhead[R]{\small   \textit{GLB-502: Igneous Petrology, Mineralogy \& Crystallography}}
\fancyfoot[C]{\small Page  \thepage\ of \pageref{LastPage}}


\newlist{parts}{enumerate}{2}
\setlist[parts,1]{label=(\alph*),leftmargin=*,itemsep=5pt,topsep=3pt}

\newlist{romanparts}{enumerate}{2}
\setlist[romanparts,1]{label=(\roman*),leftmargin=*,itemsep=5pt,topsep=3pt}

\newcommand{\pts}[1]{\hfill   \textit{[#1]}}


\begin{document}

\begin{center}
  {\large\bfseries BANARAS HINDU UNIVERSITY}\\[2pt]
  {
\normalsize B.Sc. (Hons.) --- Semester V Examination, 2013-14}\\[6pt]
  \rule{0.8\linewidth}{0.5pt}\\[6pt]
  {\large\bfseries Paper No. GLB-502: Igneous Petrology, Mineralogy and Crystallography}\\[4pt]
  {
\normalsize Time: 3 Hours \hfill Maximum Marks: 70}\\[2pt]
  {\small Ref. No.: BS/Sem V/258}
\end{center}
\rule{\linewidth}{0.4pt}
\smallskip



\noindent   \textit{Instructions: Attempt any   \textbf{five} questions, selecting at least   \textbf{two} from each of Sections A and B. Section C is compulsory. All questions carry equal marks (14 marks each).}
\bigskip

\bigskip


\noindent {\large\bfseries SECTION --- A}
\rule{\linewidth}{0.4pt}\smallskip




\noindent   \textbf{Question 1.} Give the salient features of mineralogical classification schemes for plutonic igneous rocks. \pts{14}

\medskip


\noindent   \textbf{Question 2.} Write the petrographic description of   \textbf{any four} of the following: \pts{$4  \times 3.5 = 14$}

\begin{parts}
  \item Alkali Granite
  \item Peridotite
  \item Komatiite
  \item Dacite
  \item Norite
  \item Rhyolite
\end{parts}

\medskip


\noindent   \textbf{Question 3.} Write descriptive notes on the following: \pts{$2  \times 7 = 14$}

\begin{parts}
  \item Total Alkali-Silica (TAS) classification
  \item Physical properties affecting the nature of magma
\end{parts}

\medskip


\noindent   \textbf{Question 4.} Describe with the help of a neat sketch the crystallization of magma in the phase equilibrium of Diopside-Albite-Anorthite system. Comment on its petrogenetic significance. \pts{14}

\bigskip


\noindent {\large\bfseries SECTION --- B}
\rule{\linewidth}{0.4pt}\smallskip




\noindent   \textbf{Question 5.} Describe the symmetry and forms belonging to the holosymmetric class of the Triclinic system and draw a suitable stereogram. Name any four minerals belonging to this class. \pts{14}

\medskip


\noindent   \textbf{Question 6.} Write short notes on   \textbf{any two} of the following: \pts{$2  \times 7 = 14$}

\begin{romanparts}
  \item Classification of minerals
  \item Isomorphism and Polymorphism
  \item Pyroxene group of minerals
\end{romanparts}

\medskip


\noindent   \textbf{Question 7.} Write a detailed note on the structural classification of silicates. \pts{14}

\bigskip


\noindent {\large\bfseries SECTION --- C (Compulsory)}
\rule{\linewidth}{0.4pt}\smallskip




\noindent   \textbf{Question 8.} Answer the following questions (in two or three sentences, or by filling in the blanks as indicated): \pts{$7  \times 2 = 14$}

\begin{romanparts}
  \item What is secondary magma?
  \item What is a rock series?
  \item What are the petrographic differences between syenite and alkali feldspar syenite?
  \item The most abundant element in the earth's crust is \underline{\hspace{3cm}} whereas the most abundant mineral in the Earth's crust is \underline{\hspace{3cm}}.
  \item What is relief in optical mineralogy?
  \item Name two minerals showing pleochroism.
  \item Name two biaxial minerals.
\end{romanparts}

\vfill\rule{\linewidth}{0.4pt}

\begin{center}\small
\href{https://bhu-exam.vercel.app/}{Click here to visit our website}\\[4pt]
\href{https://me-aryan.vercel.app/}{Click here to visit our contributor}\\[4pt]
\href{https://quashan-me.vercel.app/}{Click here to visit our contributor}
\end{center}
\end{document}
